% ----------------------------------------------------------
% Capítulo 5 — Resultados
% Tabelas, figuras e valores citados na prosa (\res{...}) são GERADOS por pfc-relatorio
% (pfc_busca/src/pfc_busca/evaluation/report.py) a partir de results/:
%   rodadas completas:  pfc-relatorio --modelos qwen3:4b-instruct-2507-q4_K_M,gemma4:e4b-it-qat,gemma4:e2b-it-qat,mistral-nemo:12b --comparar-com results/estacao --paper ../paper_revisado
%   estação:            pfc-relatorio --resultados results/estacao --modelos qwen3:4b-instruct-2507-q4_K_M,gemma4:e2b-it-qat --sufixo _estacao --titulo " --- estação de referência (camadas P e N)" --paper ../paper_revisado
%   nuvem:              pfc-relatorio --modelos <modelos do Groq com rodada completa> --sufixo _groq --titulo " --- referência em nuvem (Groq)" --paper ../paper_revisado
%   nuvem, mesmas consultas das rodadas completas:
%                       pfc-relatorio --modelos <idem> --sufixo _groqcomum --titulo " --- referência em nuvem, consultas das rodadas completas" --mesmas-consultas-de results --modelos-referencia <modelos locais> --saida results/consolidado_groqcomum --paper ../paper_revisado
%   reexecução na estação: ver o comentário da seção "Validação na Estação de Referência".
% Nenhum número deste capítulo é digitado à mão. A prosa interpretativa foi escrita sobre os valores
% gerados e deve ser relida sempre que as rodadas forem refeitas.
% ----------------------------------------------------------
\chapter{Resultados}
\label{cap:resultados}

% Guardas: o capítulo compila mesmo antes de os artefatos existirem.
\newcommand{\tabelagerada}[2]{\IfFileExists{tabelas/#1.tex}{\input{tabelas/#1}}{\par\noindent[DRAFT --- tabela \texttt{\detokenize{#1}} gerada por \texttt{pfc-relatorio}; #2]\par}}
\newcommand{\figuragerada}[3]{\IfFileExists{figuras/#1.pdf}{%
\begin{figure}[htbp!]\centering\includegraphics[width=0.92\textwidth]{figuras/#1.pdf}\caption{#2}\label{fig:#1}\fonte{Elaborado pelo autor.}\end{figure}}{\par\noindent[DRAFT --- figura \texttt{\detokenize{#1}}; #3]\par}}

Este capítulo apresenta os resultados da execução do protocolo descrito no Capítulo~\ref{cap:metodologia} sobre o \textit{dataset} do Apêndice~\ref{apendice:dataset}. Todos os valores são produzidos pelo \textit{pipeline} de avaliação (Seção~\ref{sec:pipeline-avaliacao}) a partir dos registros de execução, com a pontuação recalculada contra o gabarito vigente; cada rodada carrega um manifesto com a \textit{tag} e o \textit{digest} do modelo, as versões de \textit{software}, o \textit{hardware}, a data de referência e os \textit{hashes} do \textit{dataset} e da definição da ferramenta.

% ---
\section{Configuração Experimental}
\label{sec:config-experimental}
% ---

Os resultados vêm de três conjuntos de rodadas, descritos na Seção~\ref{sec:ambiente}: as \textbf{rodadas completas} dos quatro modelos locais, com três repetições, executadas com Ollama numa GPU de nuvem NVIDIA T4; a \textbf{validação na estação de referência}, com as 62 consultas das camadas P e N; e a \textbf{referência em nuvem}, com \res{groq}{geral}{nmodelos} modelos de maior porte servidos pelo Groq, com uma repetição. As Tabelas~\ref{tab:ambiente}, \ref{tab:ambiente_estacao} e~\ref{tab:ambiente_groq} transcrevem dos manifestos o ambiente de cada rodada. \res{principal}{geral}{frasecobertura}

\tabelagerada{tab_ambiente}{ambiente das rodadas completas}

Cada modelo foi avaliado com temperatura zero, limite de 1024 \textit{tokens} de saída, sem contexto de conversações anteriores (\textit{zero-shot}), sem exemplos no \textit{prompt} e sem dicionário de normalização, com o mesmo \textit{prompt} de sistema e a mesma definição da ferramenta \texttt{buscar\_catalogo}; o modo de raciocínio foi desativado nos modelos locais que o oferecem (Seção~\ref{sec:config-modelos}). As \textit{tags} exatas são \res{principal}{geral}{tags}. Nas rodadas completas, as métricas principais cobrem \res{principal}{geral}{nprincipais} consultas, pontuadas em cada uma das três repetições --- \res{principal}{qwen}{execucoes} execuções por modelo ---, incluindo as da categoria F, em que a resposta correta é não chamar a ferramenta; as consultas observacionais são reportadas à parte (Seção~\ref{sec:observacionais}). Nenhuma execução das rodadas completas terminou em erro de infraestrutura.

% ---
\section{Comparativo Geral entre Modelos}
\label{sec:comparativo}
% ---

A Tabela~\ref{tab:resultados_comparativo} sumariza as métricas principais por modelo, e a Tabela~\ref{tab:latencia} detalha a distribuição da latência da tradução.

\tabelagerada{tab_comparativo}{comparativo geral}

\tabelagerada{tab_latencia}{latência por modelo}

\paragraph{Qualidade.} Três dos quatro modelos ficaram num mesmo patamar de acurácia: \res{principal}{e4b}{acuracia} para o Gemma 4 E4B, \res{principal}{qwen}{acuracia} para o Qwen 3 4B e \res{principal}{e2b}{acuracia} para o Gemma 4 E2B, com intervalos de confiança que se sobrepõem amplamente. Pelo teste de McNemar (Seção~\ref{sec:estatistica}), nenhuma das diferenças entre esses três é estatisticamente significativa. O Mistral Nemo 12B, o maior dos quatro, ficou muito abaixo, com \res{principal}{mistral}{acuracia} (IC 95\%: \res{principal}{mistral}{ic}), e a diferença para cada um dos outros três é significativa. O tamanho do modelo, portanto, não explica o desempenho na tarefa: a Seção~\ref{sec:analise-erros} mostra que a queda do Mistral Nemo decorre quase inteiramente de um modo de falha no uso da ferramenta.

\paragraph{Precisão e \textit{recall}.} Os três modelos do patamar superior têm F1 ponderado semelhante (\res{principal}{qwen}{f1}, \res{principal}{e4b}{f1} e \res{principal}{e2b}{f1}), mas chegam a ele por caminhos diferentes. O Gemma 4 E4B tem a maior precisão (\res{principal}{e4b}{precisao}) e o menor \textit{recall} (\res{principal}{e4b}{recall}) dos três: quando chama a ferramenta, raramente preenche um campo errado, mas deixa de chamá-la com frequência. O Qwen 3 4B tem o maior \textit{recall} (\res{principal}{qwen}{recall}) e a menor precisão (\res{principal}{qwen}{precisao}), porque acrescenta com mais frequência campos que a consulta não pediu.

\paragraph{Latência.} Na T4, o Gemma 4 E2B foi o mais rápido, com mediana de \res{principal}{e2b}{latmeds}~s e percentil 95 de \res{principal}{e2b}{latpnoventaecinco}~s; o Gemma 4 E4B e o Qwen 3 4B tiveram medianas de \res{principal}{e4b}{latmeds}~s e \res{principal}{qwen}{latmeds}~s; e o Mistral Nemo 12B, \res{principal}{mistral}{latmeds}~s, cerca de \res{principal}{e2b-mistral}{razaolat} vezes a mediana do Gemma 4 E2B. A cauda do Gemma 4 E4B é mais longa que a do Qwen (percentil 95 de \res{principal}{e4b}{latpnoventaecinco}~s, contra \res{principal}{qwen}{latpnoventaecinco}~s), como mostra a Figura~\ref{fig:fig_latencia_boxplot}. Esses tempos valem para uma GPU com VRAM suficiente para cada modelo inteiro; na estação de referência, com 4~GB, a latência é maior e depende da parcela do modelo que cabe na GPU (Seção~\ref{sec:validacao-local}).

\figuragerada{fig_latencia_boxplot}{Distribuição da latência da tradução nas rodadas completas, em escala logarítmica, por modelo}{boxplot de latência}

\paragraph{Estabilidade entre repetições.} Numa mesma sessão da T4, as três repetições produziram exatamente a mesma resposta em \res{principal}{qwen}{estaveis} das \res{principal}{qwen}{nestaveis} consultas para o Qwen 3 4B, em \res{principal}{e2b}{estaveis} para o Gemma 4 E2B, em \res{principal}{mistral}{estaveis} para o Mistral Nemo 12B e em \res{principal}{e4b}{estaveis} para o Gemma 4 E4B. Com temperatura zero e o mesmo ambiente, a variação entre repetições foi, portanto, pequena; a Seção~\ref{sec:validacao-local} mostra que ela é maior entre sessões diferentes na estação de referência.

% ---
\section{Análise por Categoria de Consulta}
\label{sec:por-categoria}
% ---

A Tabela~\ref{tab:resultados_por_categoria} apresenta o F1 ponderado por categoria, e a Figura~\ref{fig:fig_acuracia_por_categoria}, a acurácia por consulta. A categoria F aparece só na figura: como o gabarito dessas consultas não tem campos, a acurácia nela é a taxa de recusa correta.

\tabelagerada{tab_por_categoria}{F1 por categoria}

\figuragerada{fig_acuracia_por_categoria}{Acurácia por consulta, por categoria e por modelo, nas rodadas completas}{acurácia por categoria}

\paragraph{O que é fácil.} As consultas simples, com um único parâmetro, tiveram as maiores acurácias entre as categorias com campos: \res{principal}{e4b}{accS} no Gemma 4 E4B e \res{principal}{qwen}{accS} no Qwen 3 4B e no Gemma 4 E2B. Os quatro modelos recusaram corretamente as consultas fora do domínio --- a acurácia na categoria F foi de \res{principal}{qwen}{accF}, \res{principal}{e4b}{accF}, \res{principal}{e2b}{accF} e \res{principal}{mistral}{accF} para Qwen 3 4B, Gemma 4 E4B, Gemma 4 E2B e Mistral Nemo 12B ---, o que confirma que o \textit{Tool Calling} permite ao modelo decidir não buscar, algo que a Saída Estruturada do protótipo não permitia.

\paragraph{O que é difícil para todos.} As categorias de tempo relativo e de ordenação foram as mais difíceis para os quatro modelos. No tempo relativo, a melhor acurácia foi a do Qwen 3 4B, \res{principal}{qwen}{accT}; na ordenação, nenhum modelo passou de \res{principal}{qwen}{accO}, e o Gemma 4 E2B não acertou nenhuma consulta dessa categoria. As consultas compostas, que exigem vários campos corretos ao mesmo tempo, ficaram em torno de \res{principal}{qwen}{accC} a \res{principal}{e4b}{accC} nos três melhores modelos. Essas quedas simultâneas indicam dificuldades da tarefa, e não de um modelo específico: resolver uma expressão temporal em datas ISO e decidir se, e por qual data, ordenar são as etapas em que a tradução mais depende de raciocínio além do reconhecimento de nomes.

\paragraph{Onde os modelos divergem.} Nas consultas com código MI ou INOM, o Qwen 3 4B foi o mais consistente (\res{principal}{qwen}{accM}, contra \res{principal}{e2b}{accM} do Gemma 4 E2B e \res{principal}{e4b}{accM} do Gemma 4 E4B). Nas consultas ambíguas e informais, o Gemma 4 E4B teve a maior acurácia (\res{principal}{e4b}{accA}).

% ---
\section{Análise por Campo do \textit{Schema}}
\label{sec:por-campo}
% ---

\tabelagerada{tab_por_campo}{F1 por campo}

\figuragerada{fig_f1_por_campo}{F1-\textit{score} por campo do \textit{schema} de busca, por modelo, nas rodadas completas}{F1 por campo}

Os campos de identificação --- \texttt{supplyArea}, \texttt{city}, \texttt{project} e \texttt{keyword} --- tiveram, em geral, os maiores F1 nos três melhores modelos: o Gemma 4 E2B acertou todas as ocorrências de \texttt{supplyArea} (F1 de \res{principal}{e2b}{f1supplyArea}), e o Qwen 3 4B teve o maior F1 em \texttt{keyword} (\res{principal}{qwen}{f1keyword}), coerente com seu desempenho nas consultas com código MI/INOM. São campos em que a tarefa se reduz a reconhecer e normalizar um nome, justamente o que o dicionário fixo do protótipo fazia; aqui, a normalização foi feita pelo próprio modelo, a partir das descrições da ferramenta.

Os campos de período e de ordenação concentraram as perdas. O F1 de \texttt{publicationPeriod} ficou entre \res{principal}{e4b}{f1publicationPeriod} (Gemma 4 E4B) e \res{principal}{e2b}{f1publicationPeriod} (Gemma 4 E2B), e os de \texttt{sortField} e \texttt{sortDirection} não passaram de \res{principal}{qwen}{f1sortField} em nenhum dos quatro modelos. O campo \texttt{creationPeriod} tem poucas ocorrências nas consultas pontuadas, o que torna seu F1 instável.

% ---
\section{Análise de Erros}
\label{sec:analise-erros}
% ---

A classificação dos erros (Seção~\ref{sec:metricas}) mostra um modo de falha dominante diferente em cada modelo. Os exemplos citados são da primeira repetição e estão, com os demais, no arquivo \texttt{results/consolidado/erros\_representativos.md} do repositório.

\paragraph{Mistral Nemo 12B: a chamada descrita em texto.} O Mistral Nemo 12B deixou de chamar a ferramenta em \res{principal}{mistral}{erronaochamou} das \res{principal}{mistral}{execucoes} execuções pontuadas. Em \res{principal}{mistral}{naochamounarrada} delas, a resposta cita a ferramenta \texttt{buscar\_catalogo} e descreve, em texto corrido, a chamada que faria. Para ``preciso da carta MI 2965-2-NE do segundo cgeo'', por exemplo, o modelo respondeu que chamaria a ferramenta com \texttt{keyword} ``2965-2-NE'' e \texttt{supplyArea} ``2° Centro de Geoinformação'', mas não emitiu a chamada estruturada. Em outras \res{principal}{mistral}{naochamouesclarecimento}, pediu mais informações ao usuário. Nesses casos, o modelo identifica a ferramenta e os parâmetros, mas não emite a chamada no canal de \textit{tool calls} da integração com o Ollama; como o protocolo não interpreta texto livre como chamada, cada um deles conta como falha.

\paragraph{Gemma 4 E4B e E2B: o pedido de esclarecimento.} Os dois modelos Gemma deixaram de chamar a ferramenta em \res{principal}{e4b}{erronaochamou} (E4B) e \res{principal}{e2b}{erronaochamou} (E2B) execuções, quase sempre pedindo ao usuário que especificasse melhor a busca --- \res{principal}{e4b}{naochamouesclarecimento} e \res{principal}{e2b}{naochamouesclarecimento} casos, respectivamente. Isso ocorreu sobretudo em consultas amplas, com um único critério: ``cartas topográficas'', ``produtos em escala 1:50.000'', ``modelos digitais de terreno''. O Gemma 4 E2B cometeu ainda um erro próprio, o campo acrescentado mais frequente nas suas respostas: copiar para \texttt{keyword} palavras comuns da consulta, como ``mdt'' em ``mdt do para em 25k'' ou ``mapeamento detalhado'' em ``mapeamento detalhado (25k ou 50k) de brasilia''.

\paragraph{Qwen 3 4B: a recusa indevida e os campos acrescentados.} O Qwen 3 4B chamou a ferramenta com mais regularidade, mas declarou fora do escopo \res{principal}{qwen}{naochamouescopo} execuções de consultas legítimas, como ``primeiro mapeamento feito em cuiaba'' e ``ultima atualizacao do para''. Seu erro mais frequente, porém, foi acrescentar campos não pedidos, sobretudo ordenação e limite (\texttt{sortField}, \texttt{sortDirection} e \texttt{limit} são os campos mais frequentes entre os seus falsos positivos): para ``todas as cartas em escala maior que 100k do rio grande do sul'', emitiu ordenação decrescente por data de publicação e limite de 100 resultados, e omitiu a escala.

\paragraph{Datas relativas.} Os erros em períodos seguem padrões comuns aos modelos: tratar ``esse ano'' como começando na data de referência, e não em 1º de janeiro; reduzir ``2 anos atrás'' a um único dia; e trocar o trimestre-calendário por três meses corridos, ou pelo trimestre errado, em ``trimestre passado'' e ``último trimestre do ano passado''. A Tabela~\ref{tab:resultados_por_campo} mostra o efeito desses erros no F1 dos campos de período.

Os tipos de erro têm, assim, remédios diferentes. A chamada descrita em texto do Mistral Nemo é uma limitação da integração do modelo com o canal de ferramentas. O pedido de esclarecimento dos modelos Gemma e a recusa indevida do Qwen são decisões de chamar ou não, sensíveis à redação do \textit{prompt}. Os erros de período e de ordenação são erros de raciocínio sobre a consulta. Nenhum deles decorre de valores fora do \textit{schema}: não houve, em nenhum modelo, argumento com campo inexistente ou tipo inválido.

\subsection{Casos observacionais}
\label{sec:observacionais}

As quatro consultas observacionais testam o comportamento diante de pedidos que a ferramenta não consegue representar. O Qwen 3 4B foi o único que chamou a ferramenta nas duas consultas com valor impossível, em \res{principal}{qwen}{obschamou} das \res{principal}{qwen}{obsn} execuções observacionais: para ``escala 1:10.000.000'', emitiu a escala 1:10.000, que existe no enumerado mas não corresponde ao pedido; para ``cartas do estado 42'', emitiu \texttt{state} ``Rio Grande do Sul'', um valor sem nenhum apoio no texto. Os demais modelos não chamaram a ferramenta nesses dois casos. Nas consultas de resposta discutível --- ``cartas topo no oceano atlântico'' e ``cartas do futuro (ano 2050)'' ---, os modelos em geral também não chamaram a ferramenta; o Gemma 4 E2B buscou cartas topográficas sem restrição de área na primeira delas.

% ---
\section{Validação na Estação de Referência}
\label{sec:validacao-local}
% ---

% Reexecução (tab_concordancia_reexecucao, numeros_reexecucao):
%   pfc-relatorio --resultados results/estacao --modelos qwen3:4b-instruct-2507-q4_K_M --comparar-com results/_descartados/estacao_14set_vram_disputada --nome-comparacao reexecucao --rotulos-comparacao "24/set,14/set" --legenda-comparacao "Qwen 3 4B executado duas vezes na estação de referência, com os mesmos pesos, prompt, ferramenta e data de referência (camadas P e N)" --so-comparacao --saida results/estacao/consolidado --paper ../paper_revisado
As 62 consultas das camadas P e N foram executadas na estação de referência pelos dois modelos de menor porte, Qwen 3 4B e Gemma 4 E2B, com uma repetição e a mesma data de referência das rodadas completas. As duas rodadas foram feitas em sequência, na mesma sessão. Antes de cada uma, o \textit{pipeline} descarregou do Ollama todos os modelos residentes, de modo que cada modelo foi carregado a frio, com a VRAM que a estação tinha livre (Tabela~\ref{tab:ambiente_estacao}). Com 4~GB de VRAM, nenhum dos dois coube inteiro na GPU: o Ollama alocou na VRAM \res{estacao}{qwen}{nagpu} do Qwen 3 4B carregado, contando pesos e contexto, e \res{estacao}{e2b}{nagpu} do Gemma 4 E2B; o restante executou na CPU.

\tabelagerada{tab_ambiente_estacao}{ambiente da validação na estação}

\tabelagerada{tab_comparativo_estacao}{comparativo na estação}

\tabelagerada{tab_latencia_estacao}{latência na estação}

\paragraph{Qualidade.} Nas \res{estacao}{qwen}{n} consultas pontuadas, os dois modelos tiveram acurácias estatisticamente indistinguíveis: \res{estacao}{qwen}{acuracia} para o Qwen 3 4B (IC 95\%: \res{estacao}{qwen}{ic}) e \res{estacao}{e2b}{acuracia} para o Gemma 4 E2B (IC 95\%: \res{estacao}{e2b}{ic}). Entre as consultas acertadas por apenas um deles, \res{estacao}{qwen-e2b}{soa} foram do Qwen e \res{estacao}{qwen-e2b}{sob} do Gemma, com \textit{p}-valor de \res{estacao}{qwen-e2b}{p} no teste de McNemar exato (Tabela~\ref{tab:mcnemar_estacao}).

\tabelagerada{tab_mcnemar_estacao}{McNemar na estação}

\paragraph{Latência.} A latência, ao contrário, separou nitidamente os dois modelos. O Qwen 3 4B teve mediana de \res{estacao}{qwen}{latmeds}~s e percentil 95 de \res{estacao}{qwen}{latpnoventaecinco}~s; o Gemma 4 E2B, \res{estacao}{e2b}{latmeds}~s e \res{estacao}{e2b}{latpnoventaecinco}~s, cerca de \res{estacao}{qwen-e2b}{razaolat} vezes mais rápido, apesar da menor parcela na GPU. Dois fatores são compatíveis com essa diferença: o Gemma 4 E2B computa com cerca de 2 bilhões de parâmetros efetivos (Seção~\ref{sec:selecao-modelos}), e suas respostas foram mais curtas, com mediana de \res{estacao}{e2b}{toksaida} \textit{tokens} de saída, contra \res{estacao}{qwen}{toksaida} do Qwen 3 4B. Em relação à T4, a latência mediana do Qwen 3 4B subiu de \res{principal}{qwen}{latmeds}~s para \res{estacao}{qwen}{latmeds}~s, e a do Gemma 4 E2B, de \res{principal}{e2b}{latmeds}~s para \res{estacao}{e2b}{latmeds}~s.

\figuragerada{fig_latencia_boxplot_estacao}{Distribuição da latência da tradução na estação de referência, em escala logarítmica}{boxplot de latência na estação}

\paragraph{Comportamento.} Na categoria F, a acurácia foi de \res{estacao}{qwen}{accF} para o Qwen 3 4B e de \res{estacao}{e2b}{accF} para o Gemma 4 E2B: os dois recusaram as consultas fora do domínio das camadas P e N. Os modos de falha foram os mesmos das rodadas completas: o Qwen declarou fora do escopo pedidos como ``atualizações deste mês'', e o Gemma pediu mais critérios em consultas como ``cartas topográficas'' e ``produtos publicados nos últimos 3 meses'', contrariando a instrução do \textit{prompt} de sistema de que um único parâmetro basta para chamar a ferramenta.

\figuragerada{fig_acuracia_por_categoria_estacao}{Acurácia por consulta, por categoria, na estação de referência}{acurácia por categoria na estação}

\paragraph{Comparação com as rodadas completas.} A Tabela~\ref{tab:concordancia-ambientes} compara, para os mesmos modelos e consultas, as respostas da estação com as da primeira repetição das rodadas completas. A resposta foi idêntica em \res{ambientes}{qwen}{identicas} das \res{ambientes}{qwen}{n} consultas para o Qwen 3 4B e em \res{ambientes}{e2b}{identicas} para o Gemma 4 E2B, e a pontuação coincidiu em \res{ambientes}{qwen}{mesmanota} e \res{ambientes}{e2b}{mesmanota}, respectivamente. A acurácia nessas consultas foi de \res{ambientes}{qwen}{acuraciaa} e \res{ambientes}{e2b}{acuraciaa} nas rodadas completas, e de \res{ambientes}{qwen}{acuraciab} e \res{ambientes}{e2b}{acuraciab} na estação. A qualidade da tradução, portanto, praticamente não mudou com o \textit{hardware}; a latência, sim.

\tabelagerada{tab_concordancia_ambientes}{concordância entre a estação e as rodadas completas}

\paragraph{Reexecução na mesma estação.} O Qwen 3 4B já havia executado as mesmas consultas na estação em 14 de setembro de 2026, com os mesmos pesos, \textit{prompt}, ferramenta, data de referência e versão do Ollama. A Tabela~\ref{tab:concordancia-reexecucao} compara as duas execuções: mesmo com temperatura zero, a resposta se repetiu em \res{reexecucao}{qwen}{identicas} das \res{reexecucao}{qwen}{n} consultas, e a pontuação, em \res{reexecucao}{qwen}{mesmanota}. A explicação mais provável são diferenças numéricas de ponto flutuante na inferência, que dependem da divisão do modelo entre GPU e CPU e da ordem das operações e bastam para mudar a escolha de um \textit{token} em respostas isoladas. A variação é maior entre sessões na estação, com parte do modelo na CPU, do que entre as repetições de uma mesma sessão na T4 (Seção~\ref{sec:comparativo}), e é por causa dela que as rodadas completas usam três repetições e que a acurácia é reportada com intervalo de confiança.

\tabelagerada{tab_concordancia_reexecucao}{reexecução do Qwen 3 4B na estação}

% ---
\section{Referência em Nuvem}
\label{sec:referencia-nuvem}
% ---

\res{groq}{geral}{nmodelosmaiusc} modelos de maior porte servidos pelo Groq executaram o mesmo protocolo, com uma repetição: \res{groq}{geral}{tags}. Eles não atendem ao requisito de execução local (RNF1) e servem como referência de quanto modelos maiores, sem restrição de \textit{hardware}, acertam nas mesmas consultas; a latência inclui a rede e não é comparável à dos modelos locais.

\tabelagerada{tab_ambiente_groq}{ambiente da referência em nuvem}

\tabelagerada{tab_comparativo_groq}{comparativo da referência em nuvem}

\tabelagerada{tab_por_categoria_groq}{F1 por categoria na referência em nuvem}

Sobre todas as \res{groq}{gptoss20}{n} consultas das métricas principais, o GPT-OSS 20B acertou \res{groq}{gptoss20}{acuracia} (IC 95\%: \res{groq}{gptoss20}{ic}) e o GPT-OSS 120B, \res{groq}{gptoss120}{acuracia} (IC 95\%: \res{groq}{gptoss120}{ic}). Restrita às mesmas \res{groqcomum}{gptoss20}{n} consultas pontuadas nas rodadas completas dos modelos locais, a acurácia foi de \res{groqcomum}{gptoss20}{acuracia} e \res{groqcomum}{gptoss120}{acuracia}, respectivamente, contra \res{principal}{e4b}{acuracia} do melhor modelo local. A distância é maior justamente nas categorias difíceis para os modelos locais: nas consultas compostas e de tempo relativo desse mesmo conjunto, o GPT-OSS 20B acertou \res{groqcomum}{gptoss20}{accC} e \res{groqcomum}{gptoss20}{accT}, contra no máximo \res{principal}{e4b}{accC} e \res{principal}{qwen}{accT} entre os modelos locais. A ordenação continuou sendo a categoria mais difícil também para eles: sobre todo o \textit{dataset}, o GPT-OSS 20B acertou \res{groq}{gptoss20}{accO} dessas consultas e o GPT-OSS 120B, \res{groq}{gptoss120}{accO}. Como nas rodadas locais, todas as consultas fora do domínio foram recusadas.

A referência mostra que a tarefa, com a mesma ferramenta, o mesmo \textit{prompt} e o mesmo gabarito, é resolvível com alta acurácia por modelos de maior porte. A distância para os modelos locais mede, assim, a limitação dos modelos que cabem em estações modestas, e não um defeito do protocolo ou do gabarito.

% ---
\section{Análise Estatística}
\label{sec:estatistica}
% ---

A Tabela~\ref{tab:mcnemar} traz o teste de McNemar exato entre cada par de modelos das rodadas completas, sobre as mesmas consultas e repetições. Entre os três modelos do patamar superior, nenhuma diferença é significativa: o \textit{p}-valor foi de \res{principal}{qwen-e4b}{p} entre Qwen 3 4B e Gemma 4 E4B, de \res{principal}{qwen-e2b}{p} entre Qwen 3 4B e Gemma 4 E2B e de \res{principal}{e4b-e2b}{p} entre os dois Gemma. Entre o Mistral Nemo 12B e cada um dos outros três, o \textit{p}-valor foi \res{principal}{qwen-mistral}{p}. Como as três repetições de cada consulta entram como pares, e elas são quase sempre idênticas numa mesma sessão, os testes tratam como independentes observações que não o são inteiramente; os \textit{p}-valores devem ser lidos como aproximação, o que não altera as conclusões, dada a distância entre os dois grupos de modelos e a proximidade dentro do grupo superior.

\tabelagerada{tab_mcnemar}{McNemar entre pares de modelos}

% ---
\section{Discussão}
\label{sec:discussao}
% ---

\paragraph{Qualidade e latência.} Os resultados não apontam um modelo dominante em qualidade entre Qwen 3 4B, Gemma 4 E4B e Gemma 4 E2B: as acurácias são estatisticamente indistinguíveis, e cada um tem pontos fortes próprios --- o Qwen nos códigos MI/INOM e no \textit{recall}, o Gemma 4 E4B na precisão e nas consultas informais, o Gemma 4 E2B nos campos de escala, tipo de produto e área de suprimento. A latência os separa com mais clareza: o Gemma 4 E2B foi o mais rápido tanto na T4 quanto na estação de referência, onde sua mediana ficou abaixo de um segundo.

\paragraph{O que o \textit{Tool Calling} entregou.} Em relação à Saída Estruturada do protótipo, três ganhos aparecem diretamente nos dados. O modelo passou a poder recusar: todas as consultas fora do domínio foram recusadas por todos os modelos. A normalização de siglas, grafias sem acento, apelidos de escala e nomes de projeto passou a ser feita pelo modelo, sem dicionário, com F1 elevado nos campos de identificação. E toda falha passou a ser registrada e classificada, o que permitiu identificar o modo de falha de cada modelo --- inclusive o do Mistral Nemo 12B, que um mecanismo de \textit{fallback} teria mascarado.

\paragraph{O que continua difícil.} Datas relativas, ordenação e consultas com muitos campos permanecem difíceis para todos os modelos locais, e a ordenação o é também para os modelos em nuvem. Parte da dificuldade é inerente à tarefa: nem sempre a consulta diz se a ordem é por publicação ou por criação, e o gabarito aceita as duas leituras, mas os modelos ainda assim acrescentam ou omitem campos. Outra parte é sensível ao \textit{prompt}: a recusa indevida e o pedido de esclarecimento em consultas amplas contrariam uma instrução explícita e são candidatos naturais a ajuste de \textit{prompt} ou a exemplos no \textit{prompt}, que este trabalho não usou por desenho.

\paragraph{\textit{Hardware}.} Com 4~GB de VRAM, nenhum modelo coube inteiro na GPU da estação, e a latência depende diretamente da parcela alocada. A qualidade, por outro lado, praticamente não mudou entre a estação e a T4. Para a DSG, isso significa que a escolha do modelo pode ser feita pela qualidade medida em qualquer ambiente, e que o \textit{hardware} de implantação define a latência: uma GPU com 6 a 8~GB de VRAM já comportaria inteiros o Qwen 3 4B e o Gemma 4 E2B.

\subsection{Recomendação de Modelo para Uso Operacional}

Para o uso operacional sobre o acervo da DSG, em estações com GPU modesta, recomenda-se o \textbf{Gemma 4 E2B}: sua acurácia é estatisticamente equivalente à dos melhores modelos avaliados, e ele teve a menor latência nos dois ambientes, inclusive na estação de referência de 4~GB. O \textbf{Qwen 3 4B} é a alternativa quando as consultas com códigos MI/INOM forem predominantes, categoria em que teve o melhor desempenho, ao custo de uma latência maior na estação. O Gemma 4 E4B, de qualidade equivalente, exige mais VRAM que os outros dois e não foi medido na estação de referência. O Mistral Nemo 12B não é recomendado na configuração avaliada, pela falha sistemática no uso do canal de ferramentas.

A recomendação tem três ressalvas. Ela vale para a configuração avaliada --- sem raciocínio, sem exemplos no \textit{prompt} e com as versões registradas nos manifestos ---, e o desempenho com raciocínio ativado não foi medido. A distância para os modelos em nuvem indica margem de melhoria considerável, que pode vir de novos modelos abertos: o \textit{pipeline} permite repetir a comparação com o mesmo critério. E os modos de falha identificados, sobretudo o pedido de esclarecimento em consultas amplas, devem ser observados no uso real antes da implantação em larga escala.
